
\chapter{Цифровой рубль}\label{ch:cbdc}
Цифровой рубль -- третья форма денег наряду с~наличными и безналичными
(см.\ раздел~\ref{sec:banks-forms-of-money}): цифровая форма
национальной валюты, эмитированная Банком России
непосредственно~\cite{fz-340,cbr-digital-ruble}, с~отдельным реестром
обязательств на платформе Банка России.
Сравнение с~СБП подробно разобрано в~разделе~\ref{sec:sbp-cbdc}.

\section{Двухуровневая архитектура}\label{sec:cbdc-architecture}

В~действующей розничной модели архитектура цифрового рубля
двухуровневая. Банк России эмитирует цифровой рубль и управляет
платформой, а~банки дают клиентам доступ к~кошелькам и исполняют
их поручения на платформе Банка России. Клиент управляет кошельком
через мобильное приложение банка, подключённого к~платформе~\cite{cbr-digital-ruble}.
Требования к~защите информации для банков, участников платформы,
установлены Положением Банка России
833-П~\cite{cbr-833p}.

Отличие от безналичного рубля -- в~природе обязательства. Безналичный
рубль -- обязательство коммерческого банка перед клиентом; цифровой
рубль -- прямое обязательство Банка России. Для клиента
исчезает кредитный риск банка; для банков возникает вопрос
о~судьбе депозитной базы~\cite{cbr-digital-ruble}.

\section{Смарт-контракты}\label{sec:cbdc-smart-contracts}

Смарт-контракты реализованы в~двух режимах с~разной степенью готовности.
Простые смарт-контракты уже работают в~пилотном режиме: они автоматически
переводят цифровые рубли в~заданные дату и~время.
К~маю 2025 года в~пилоте исполнено более 17~тыс.\ таких контрактов~\cite{cbr-digital-ruble}; типовые сценарии: автоплатёж за
жилищно-коммунальные услуги, автопополнение кошелька, поэтапная оплата
подрядчику~\cite{gazprombank-smart-contracts-2025}.

Коммерческая платформа смарт-контрактов, на которой банки смогут
публиковать собственные сценарии с~конфиденциальными условиями и
юридической обёрткой, находится в~стадии разработки. В~марте 2026 года
глава Банка России сообщила, что регулятор \enquote{представит для
обсуждения концепцию платформы коммерческих смарт-контрактов} в~течение
2026 года~\cite{cbr-nabiullina-abr-2026}.
Потенциальные применения: целевые выплаты, страховые сценарии, B2B
расчёты с~условиями. Пока они описаны на уровне концепции. Офлайн-режим
также остаётся перспективной возможностью, а~не запущенным функционалом~\cite{cbr-digital-ruble}.

\section{Цифровой рубль и СБП: разные слои}\label{sec:cbdc-vs-sbp}

Цифровой рубль нередко ставят рядом с~СБП: оба сюжета
связаны с~публичной инфраструктурой и оба обещают более дешёвый
платёж. На архитектурном уровне это разные вещи. СБП переводит
безналичные рубли между банковскими счетами; цифровой рубль
живёт в~отдельном кошельке как самостоятельная форма денег. Смешивать их
в~одну модель нельзя~\cite{cbr-digital-ruble}.
